\chapter{Testare și validare}
\pagestyle{fancy}

\section{Strategia de testare}

Scopul acestui capitol este de a demonstra că sistemul implementat îndeplinește cerințele funcționale și tehnice definite anterior. Validarea nu se limitează la observarea interfeței grafice, ci urmărește verificarea controlată a componentelor software, a interfețelor API, a fluxului de date IoT și a modulului predictiv bazat pe ARIMA.

Strategia de testare utilizată este structurată pe mai multe niveluri. La baza acesteia se află verificările statice și testele unitare, care au rolul de a identifica erori de implementare în module izolate. Nivelul intermediar este reprezentat de testele de integrare, prin care sunt validate comunicația dintre API, bazele de date și serviciul de predicție. Nivelul superior este reprezentat de testele end-to-end și de validarea manuală asistată, unde este verificat fluxul complet de la simulatorul LoRaWAN până la dashboard-ul aplicației.

Această abordare este adecvată pentru proiect deoarece sistemul este compus din servicii eterogene: aplicație web Next.js, worker MQTT, baze de date PostgreSQL și InfluxDB, infrastructură ChirpStack, simulator LoRaWAN și serviciu Python pentru predicție. Prin urmare, o singură metodă de testare nu ar fi suficientă pentru a valida corect întregul sistem.

\begin{table}[htbp]
\centering
\small
\begin{tabular}{|p{0.24\textwidth}|p{0.34\textwidth}|p{0.30\textwidth}|}
\hline
\textbf{Nivel de testare} & \textbf{Scop} & \textbf{Instrumente utilizate} \\
\hline
Verificări statice & Identificarea erorilor de sintaxă, tipare și build & ESLint, TypeScript, Next.js build \\
\hline
Teste unitare & Verificarea funcțiilor izolate & Vitest, pytest \\
\hline
Teste de integrare & Validarea API-urilor și a comunicației dintre servicii & Playwright API, Docker Compose \\
\hline
Teste end-to-end & Validarea fluxurilor reale de utilizare & Playwright, browser, date demo \\
\hline
Validare IoT & Verificarea traseului simulator--MQTT--InfluxDB--dashboard & Docker logs, InfluxDB, dashboard \\
\hline
Validare predictivă & Verificarea comportamentului ARIMA & pytest, endpoint forecast, dashboard Meter \\
\hline
\end{tabular}
\caption{Strategia generală de testare}
\label{tab:strategie-testare}
\end{table}

\section{Mediul de testare}

Testarea este realizată într-un mediu local, orchestrat prin Docker Compose. Acest mediu reproduce componentele necesare pentru demonstrarea funcțională a sistemului fără a necesita contoare fizice sau gateway-uri LoRaWAN reale. Aplicația web și worker-ul MQTT sunt pornite local din directorul \texttt{web-app}, iar serviciile de infrastructură sunt pornite din directorul rădăcină al proiectului.

\begin{table}[htbp]
\centering
\small
\begin{tabular}{|p{0.28\textwidth}|p{0.58\textwidth}|}
\hline
\textbf{Componentă} & \textbf{Rol în procesul de testare} \\
\hline
Docker Compose & Pornește serviciile necesare: PostgreSQL, InfluxDB, Redis, Mosquitto, ChirpStack, Gateway Bridge, simulatorul LoRaWAN și serviciul ARIMA. \\
\hline
PostgreSQL & Stochează utilizatorii, dispozitivele, tarifele și relațiile de ownership. \\
\hline
InfluxDB & Stochează telemetria istorică și datele utilizate pentru grafice și predicții. \\
\hline
ChirpStack & Gestionează dispozitivele LoRaWAN simulate și publică evenimentele uplink. \\
\hline
Mosquitto MQTT & Permite transmiterea asincronă a mesajelor de telemetrie către worker. \\
\hline
LWN Simulator & Generează trafic LoRaWAN simulat pentru contoare de electricitate, gaz, apă, încălzire și răcire. \\
\hline
Next.js Web App & Expune API-urile și interfața dashboard-ului. \\
\hline
Forecast Service & Generează predicții pe termen scurt folosind modelul ARIMA. \\
\hline
\end{tabular}
\caption{Componentele mediului local de testare}
\label{tab:mediu-testare}
\end{table}

Pregătirea mediului se realizează prin următoarele comenzi:

\begin{verbatim}
docker compose up -d
cd web-app
npm run provision:demo-devices
npm run seed:history
npm run dev:all
\end{verbatim}

Aceste comenzi pornesc infrastructura locală, creează dispozitivele demonstrative, generează date istorice și rulează simultan aplicația web și worker-ul MQTT.

Starea serviciilor pornite prin Docker Compose este utilizată ca dovadă pentru disponibilitatea mediului de testare. Figura \ref{fig:testare-docker-compose} confirmă faptul că serviciile principale ale arhitecturii au fost active în momentul validării.

\begin{figure}[htbp]
\centering
\includegraphics[width=0.95\textwidth,height=0.75\textheight,keepaspectratio]{figures/capitolul6_testare_validare/01_docker_compose_ps.png}
\caption{Starea serviciilor Docker Compose utilizate în testare}
\label{fig:testare-docker-compose}
\end{figure}

\section{Testarea componentelor software}

Testarea componentelor software urmărește verificarea modulelor care pot fi validate izolat, fără dependență directă de bazele de date sau de infrastructura LoRaWAN. Pentru aplicația web au fost introduse teste unitare cu Vitest, iar pentru serviciul de predicție sunt utilizate teste Python cu pytest.

\subsection{Verificări statice și de build}

Pentru aplicația Next.js sunt utilizate trei verificări principale:

\begin{verbatim}
npm run lint
npx tsc --noEmit
npm run build
\end{verbatim}

Comanda \texttt{npm run lint} validează respectarea regulilor ESLint. Comanda \texttt{npx tsc --noEmit} verifică tipurile TypeScript fără a genera fișiere de ieșire. Comanda \texttt{npm run build} validează faptul că aplicația poate fi compilată în mod de producție.

\subsection{Teste unitare pentru aplicația web}

Testele unitare acoperă funcții de normalizare, validare și formatare. Acestea sunt importante deoarece erorile în astfel de funcții pot afecta mai multe zone ale aplicației: API-uri, formulare, grafice și calcule de cost.

\begin{table}[htbp]
\centering
\small
\begin{tabular}{|p{0.30\textwidth}|p{0.46\textwidth}|p{0.14\textwidth}|}
\hline
\textbf{Funcționalitate testată} & \textbf{Criteriu de acceptare} & \textbf{Instrument} \\
\hline
Tipuri de utilități & Sunt acceptate doar tipurile cunoscute, iar unitățile implicite sunt corecte. & Vitest \\
\hline
Coduri de revendicare & Codurile sunt normalizate și validate înainte de utilizare. & Vitest \\
\hline
Identificator \texttt{devEui} & Formatul este normalizat la 16 caractere hexazecimale. & Vitest \\
\hline
Interogări telemetry & Intervalele de timp invalide și parametrii incompatibili sunt respinși. & Vitest \\
\hline
Formatare dashboard & Valorile monetare, cantitățile și statusurile sunt afișate defensiv. & Vitest \\
\hline
\end{tabular}
\caption{Teste unitare pentru aplicația web}
\label{tab:teste-unitare-web}
\end{table}

Rularea testelor unitare se face prin:

\begin{verbatim}
cd web-app
npm run test
\end{verbatim}

\subsection{Teste unitare pentru serviciul ARIMA}

Serviciul de predicție este testat cu pytest. Testele verifică atât cazul normal, în care există suficiente date istorice, cât și cazurile de fallback sau input invalid. Această verificare este importantă deoarece modulul ARIMA trebuie să se comporte predictibil și atunci când datele sunt insuficiente.

\begin{table}[htbp]
\centering
\small
\begin{tabular}{|p{0.34\textwidth}|p{0.42\textwidth}|p{0.14\textwidth}|}
\hline
\textbf{Scenariu} & \textbf{Rezultat așteptat} & \textbf{Instrument} \\
\hline
Serie istorică scurtă & Răspuns cu status \texttt{insufficient\_data}. & pytest \\
\hline
Serie istorică regulată & Sunt generate puncte de predicție. & pytest \\
\hline
Parametri invalizi & Serviciul returnează eroare de validare. & pytest \\
\hline
Orizont parțial & Numărul de pași este calculat prin rotunjire superioară. & pytest \\
\hline
Selecție ordin ARIMA & Ordinul selectat rămâne în grila configurată. & pytest \\
\hline
\end{tabular}
\caption{Teste unitare pentru serviciul ARIMA}
\label{tab:teste-unitare-arima}
\end{table}

Rularea testelor se face prin:

\begin{verbatim}
cd forecast-service
python -m pytest
\end{verbatim}

\section{Testarea API-ului aplicației}

Testarea API-ului urmărește validarea contractelor expuse de aplicația Next.js. Sunt verificate atât cazuri de succes, cât și cazuri de acces neautorizat. Pentru aceste teste este necesar ca aplicația să fie pornită local și baza de date să conțină utilizatorii și dispozitivele demonstrative.

\begin{table}[htbp]
\centering
\small
\begin{tabular}{|p{0.31\textwidth}|p{0.34\textwidth}|p{0.23\textwidth}|}
\hline
\textbf{Endpoint} & \textbf{Scopul testului} & \textbf{Rezultat așteptat} \\
\hline
\texttt{POST /api/auth/login} & Autentificarea contului demo de companie. & Status 200 și sesiune validă. \\
\hline
\texttt{GET /api/auth/me} & Verificarea utilizatorului autentificat. & Email-ul corespunde contului demo. \\
\hline
\texttt{GET /api/devices} & Listarea dispozitivelor asociate contului. & Lista conține dispozitive. \\
\hline
\texttt{GET /api/devices/summary} & Obținerea sumarului de flotă. & Totalul de dispozitive este pozitiv. \\
\hline
\texttt{GET /api/devices/[devEui]/readings} & Citirea telemetriei pentru un dispozitiv. & Răspuns valid pentru dispozitivul selectat. \\
\hline
\texttt{GET /api/devices/[devEui]/cost} & Calculul costului estimat. & Cost numeric mai mare sau egal cu zero. \\
\hline
\texttt{GET /api/devices/[devEui]/forecast} & Obținerea predicției ARIMA. & Status controlat: \texttt{ok}, \texttt{insufficient\_data}, \texttt{model\_error} sau \texttt{service\_unavailable}. \\
\hline
Endpoint protejat fără sesiune & Verificarea autentificării obligatorii. & Status 401. \\
\hline
Dispozitiv al altui utilizator & Verificarea izolării ownership-ului. & Acces interzis sau dispozitiv ascuns. \\
\hline
\end{tabular}
\caption{Scenarii de testare pentru API}
\label{tab:teste-api}
\end{table}

Rularea testelor API se face prin:

\begin{verbatim}
cd web-app
npm run test:e2e -- --project=api
\end{verbatim}

\section{Testarea fluxului IoT}

Validarea fluxului IoT este esențială deoarece sistemul nu este doar o aplicație web, ci o platformă care integrează o infrastructură LoRaWAN simulată. Fluxul verificat este:

\begin{verbatim}
LWN Simulator -> ChirpStack Gateway Bridge -> ChirpStack -> Mosquitto MQTT
-> MQTT Worker -> InfluxDB -> Next.js API -> Dashboard
\end{verbatim}

Pentru validare se verifică succesiv fiecare punct al traseului. Mai întâi se confirmă că serviciile Docker sunt pornite. Apoi se verifică existența dispozitivelor în ChirpStack și publicarea mesajelor uplink prin MQTT. În continuare, se verifică logurile worker-ului și existența măsurătorii \texttt{meter\_reading} în InfluxDB. La final, se verifică afișarea datelor în dashboard.

\begin{table}[htbp]
\centering
\small
\begin{tabular}{|p{0.28\textwidth}|p{0.40\textwidth}|p{0.20\textwidth}|}
\hline
\textbf{Punct verificat} & \textbf{Metodă de verificare} & \textbf{Criteriu de acceptare} \\
\hline
Containere Docker & \texttt{docker compose ps} & Serviciile sunt în stare running. \\
\hline
Dispozitive ChirpStack & Interfața ChirpStack sau scriptul de provisioning & Dispozitivele demo există. \\
\hline
Mesaje MQTT & Loguri Mosquitto/worker & Mesaje uplink recepționate. \\
\hline
Persistență InfluxDB & Interogare pentru \texttt{meter\_reading} & Există puncte pentru \texttt{devEui}. \\
\hline
Dashboard & Interfața Overview/Meter & Valorile live și istorice sunt afișate. \\
\hline
\end{tabular}
\caption{Scenarii de validare a fluxului IoT}
\label{tab:teste-flux-iot}
\end{table}

Figurile \ref{fig:testare-worker-mqtt} și \ref{fig:testare-influxdb} prezintă două dovezi complementare pentru traseul de date. Prima evidențiază conectarea worker-ului la brokerul MQTT, abonarea la topicul de uplink și apelurile API realizate de dashboard. A doua confirmă existența citirilor \texttt{meter\_reading} în InfluxDB.

\begin{figure}[htbp]
\centering
\includegraphics[width=0.95\textwidth,height=0.78\textheight,keepaspectratio]{figures/capitolul6_testare_validare/11_mqtt_worker_logs.png}
\caption{Fragment de log pentru worker-ul MQTT și apelurile API ale dashboard-ului}
\label{fig:testare-worker-mqtt}
\end{figure}

\begin{figure}[htbp]
\centering
\includegraphics[width=0.95\textwidth,height=0.78\textheight,keepaspectratio]{figures/capitolul6_testare_validare/12_influx_meter_reading.png}
\caption{Interogare InfluxDB pentru măsurătoarea \texttt{meter\_reading}}
\label{fig:testare-influxdb}
\end{figure}

\section{Testarea interfeței web}

Interfața web este validată prin teste Playwright și prin inspecție vizuală. Testele automate verifică autentificarea și accesarea rutelor principale. Inspecția vizuală este necesară deoarece dashboard-ul conține grafice, hărți și tabele a căror lizibilitate trebuie verificată în contexte diferite.

\begin{table}[htbp]
\centering
\small
\begin{tabular}{|p{0.26\textwidth}|p{0.44\textwidth}|p{0.18\textwidth}|}
\hline
\textbf{Zonă UI} & \textbf{Scenariu validat} & \textbf{Criteriu de acceptare} \\
\hline
Autentificare & Login cu contul demo de companie. & Utilizatorul ajunge în dashboard. \\
\hline
Overview & Afișarea KPI-urilor, graficelor și sumarului live. & Datele sunt vizibile și coerente. \\
\hline
Devices & Paginare, filtrare, claim panel și edit modal. & Lista este utilizabilă pentru flotă mare. \\
\hline
Meter & Selectare dispozitiv, profil consum și forecast. & Citirile și predicția sunt afișate fără blocarea paginii. \\
\hline
Billing & Costuri pe utilități și top dispozitive. & Costurile sunt afișate agregat și pe dispozitiv. \\
\hline
Map & Poziționare geografică și selecție dispozitiv. & Marcatorii sunt vizibili și selectabili. \\
\hline
Mobile & Viewport \texttt{390x844}. & Nu există clipping orizontal. \\
\hline
\end{tabular}
\caption{Scenarii de testare a interfeței web}
\label{tab:teste-ui}
\end{table}

Validarea vizuală a interfeței a fost realizată prin capturi de ecran obținute dintr-o sesiune autentificată cu utilizatorul demonstrativ de companie. Figurile \ref{fig:testare-overview}--\ref{fig:testare-mobile} prezintă principalele ecrane ale aplicației: sumarul flotei, inventarul dispozitivelor, pagina contorului individual, facturarea, vizualizarea pe hartă și afișarea pe un viewport mobil.

\begin{figure}[htbp]
\centering
\includegraphics[width=0.95\textwidth]{figures/capitolul6_testare_validare/03_overview_dashboard.png}
\caption{Validarea paginii de prezentare generală a flotei}
\label{fig:testare-overview}
\end{figure}

\begin{figure}[htbp]
\centering
\includegraphics[width=0.95\textwidth]{figures/capitolul6_testare_validare/04_devices_pagination.png}
\caption{Validarea inventarului de dispozitive și a paginării}
\label{fig:testare-devices}
\end{figure}

\begin{figure}[htbp]
\centering
\includegraphics[width=0.95\textwidth]{figures/capitolul6_testare_validare/05_meter_forecast.png}
\caption{Validarea paginii contorului individual și a prognozei ARIMA}
\label{fig:testare-meter-forecast}
\end{figure}

\begin{figure}[htbp]
\centering
\includegraphics[width=0.95\textwidth]{figures/capitolul6_testare_validare/06_billing_summary.png}
\caption{Validarea paginii de facturare și a costurilor agregate}
\label{fig:testare-billing}
\end{figure}

\begin{figure}[htbp]
\centering
\includegraphics[width=0.95\textwidth]{figures/capitolul6_testare_validare/07_map_view.png}
\caption{Validarea vizualizării geografice a dispozitivelor}
\label{fig:testare-map}
\end{figure}

\begin{figure}[htbp]
\centering
\includegraphics[width=0.45\textwidth]{figures/capitolul6_testare_validare/08_mobile_dashboard.png}
\caption{Validarea interfeței pe viewport mobil}
\label{fig:testare-mobile}
\end{figure}

Rularea testelor de interfață se face prin:

\begin{verbatim}
cd web-app
npm run test:e2e -- --project=desktop
npm run test:e2e -- --project=mobile
\end{verbatim}

\section{Validarea modulului ARIMA}

Validarea modulului ARIMA are două obiective. Primul obiectiv este verificarea funcționării serviciului Python independent. Al doilea obiectiv este verificarea integrării acestuia cu aplicația Next.js și cu interfața Meter.

Serviciul Python primește puncte istorice sub forma unor perechi timestamp-valoare. Dacă numărul de puncte valide este mai mic decât pragul minim, serviciul returnează \texttt{insufficient\_data}. Dacă există suficiente puncte, serviciul selectează ordinul ARIMA prin criteriul AIC și generează predicțiile pentru orizontul cerut.

\begin{table}[htbp]
\centering
\small
\begin{tabular}{|p{0.30\textwidth}|p{0.42\textwidth}|p{0.16\textwidth}|}
\hline
\textbf{Caz validat} & \textbf{Rezultat așteptat} & \textbf{Nivel} \\
\hline
Serie istorică suficientă & Predicție cu status \texttt{ok}. & Serviciu + API \\
\hline
Serie istorică insuficientă & Răspuns controlat \texttt{insufficient\_data}. & Serviciu + UI \\
\hline
Serviciu forecast indisponibil & Meter afișează mesaj neblocant. & API + UI \\
\hline
Cost estimat forecast & Costul este calculat cu tariful dispozitivului. & API \\
\hline
\end{tabular}
\caption{Scenarii de validare pentru modulul ARIMA}
\label{tab:validare-arima}
\end{table}

Un răspuns valid al endpoint-ului de forecast conține statusul modelului, ordinul ARIMA selectat, punctele prognozate și estimarea de cost. Aceste elemente permit atât validarea tehnică a serviciului, cât și integrarea rezultatelor în dashboard.

Disponibilitatea serviciului predictiv a fost verificată prin endpoint-ul de sănătate, prezentat în Figura \ref{fig:testare-forecast-health}. În plus, Figura \ref{fig:testare-meter-forecast} arată integrarea rezultatului în pagina contorului individual.

\begin{figure}[htbp]
\centering
\includegraphics[width=0.75\textwidth]{figures/capitolul6_testare_validare/02_forecast_health.png}
\caption{Validarea endpoint-ului de sănătate al serviciului ARIMA}
\label{fig:testare-forecast-health}
\end{figure}

\section{Rezultatele testării}

Rezultatele testării sunt centralizate în Tabelul \ref{tab:rezultate-testare}. Testele au fost rulate pe mediul local descris anterior, cu serviciile Docker pornite, date demonstrative provisionate și date istorice încărcate în InfluxDB.

\begin{table}[htbp]
\centering
\small
\begin{tabular}{|p{0.30\textwidth}|p{0.25\textwidth}|p{0.33\textwidth}|}
\hline
\textbf{Categorie test} & \textbf{Status} & \textbf{Observații} \\
\hline
Lint și TypeScript & Trecut & \texttt{npm run lint} și \texttt{npx tsc --noEmit} au rulat fără erori. \\
\hline
Build Next.js & Trecut & \texttt{npm run build} a compilat aplicația și rutele Next.js. \\
\hline
Teste unitare web-app & Trecut & Vitest a raportat 14 teste trecute. \\
\hline
Teste forecast-service & Trecut & Pytest a raportat 6 teste trecute pentru serviciul ARIMA. \\
\hline
Teste API & Trecut & Playwright API a raportat 3 teste trecute, inclusiv acces neautorizat și ownership. \\
\hline
Teste UI desktop/mobile & Trecut & Proiectul desktop a raportat 1 test trecut și 1 test sărit intenționat; proiectul mobile a raportat 1 test trecut și 1 test sărit intenționat. \\
\hline
Flux IoT end-to-end & Trecut & Serviciile Docker au fost active, worker-ul MQTT s-a conectat la broker, iar InfluxDB a returnat citiri \texttt{meter\_reading}. \\
\hline
\end{tabular}
\caption{Centralizarea rezultatelor testării}
\label{tab:rezultate-testare}
\end{table}

Dovezile de rulare pentru testele automate sunt prezentate în Figurile \ref{fig:testare-unitare} și \ref{fig:testare-api}. Acestea sunt incluse pentru trasabilitatea rezultatelor, în timp ce interpretarea finală este sintetizată în Tabelul \ref{tab:rezultate-testare}.

\begin{figure}[htbp]
\centering
\includegraphics[width=0.95\textwidth,height=0.75\textheight,keepaspectratio]{figures/capitolul6_testare_validare/10_vitest_pytest_results.png}
\caption{Rezultatele testelor unitare pentru aplicația web și serviciul ARIMA}
\label{fig:testare-unitare}
\end{figure}

\begin{figure}[htbp]
\centering
\includegraphics[width=0.95\textwidth,height=0.75\textheight,keepaspectratio]{figures/capitolul6_testare_validare/09_playwright_api_results.png}
\caption{Rezultatele testelor de integrare API}
\label{fig:testare-api}
\end{figure}

\section{Limitări ale procesului de testare}

Procesul de testare are câteva limitări. În primul rând, dispozitivele sunt simulate, nu fizice. Această alegere permite reproducerea mediului local, dar nu validează propagarea radio reală, interferențele sau comportamentul hardware. În al doilea rând, tarifele sunt statice și au rol demonstrativ, deci validarea costurilor nu reprezintă o verificare de facturare comercială. În al treilea rând, testarea de performanță este limitată la dimensiunea flotei demonstrative și nu acoperă scenarii cu mii de dispozitive.

Cu toate acestea, limitările sunt acceptabile pentru scopul lucrării, deoarece obiectivul principal este validarea unei arhitecturi software funcționale și extensibile. Sistemul demonstrează integrarea completă între simulare LoRaWAN, procesare MQTT, stocare time-series, dashboard operațional și predicție de consum.

\section{Concluzii privind validarea}

Capitolul a prezentat strategia de testare și validare a sistemului. Metodele propuse acoperă atât componentele izolate, cât și integrarea completă a platformei. Testele statice și unitare oferă încredere în corectitudinea modulelor software, testele API validează contractele aplicației, iar testele end-to-end demonstrează funcționarea scenariilor reale de utilizare.

Validarea fluxului IoT confirmă faptul că datele generate de simulator pot fi transmise prin ChirpStack și MQTT, procesate de worker, stocate în InfluxDB și afișate în dashboard. Validarea modulului ARIMA confirmă faptul că aplicația depășește nivelul de monitorizare pasivă și include o componentă de analiză predictivă.

Prin această strategie, sistemul este verificat din punct de vedere funcțional, arhitectural și operațional, iar rezultatele obținute pot fi utilizate ca suport experimental pentru concluziile lucrării.
